\documentclass[12pt]{article}
\usepackage[T1]{fontenc}
\usepackage[utf8]{inputenc}
\usepackage{lmodern}
\usepackage{textcomp}
\usepackage{enumitem}
\usepackage{geometry}
\usepackage{graphicx}
\usepackage{amsmath, amssymb}
\geometry{margin=1in}
\begin{document}
\begin{center}
Constitutional Law - Graduate Level Assessment 9 \
Total Marks: 40 \
Answer all questions.
\end{center}

\section*{Multiple Choice (10 marks)}
\begin{enumerate}[label=\arabic*.]
\item There is a grand jury proceeding underway for a local businessman. What is not true about grand jury proceedings?
\begin{enumerate}[label=(\alph*)]
    \item The proceedings are conducted in secret.
    \item There is no right to counsel.
    \item There is a right to Miranda warnings.
    \item There is no right to have evidence excluded.
\end{enumerate}
\item Which of the following constitutional provisions is applicable to corporations?
\begin{enumerate}[label=(\alph*)]
    \item The privileges and immunities clause of the Fourteenth Amendment.
    \item The comity clause of Article IV.
    \item The Fifth Amendment's prohibition against compulsory self-incrimination.
    \item The equal protection clause of the Fourteenth Amendment.
\end{enumerate}
\item A merchant's irrevocable written offer (firm offer) to sell goods
\begin{enumerate}[label=(\alph*)]
    \item Must be separately signed if the offeree supplies a form contract containing the offer.
    \item Is valid for three months.
    \item Is nonassignable.
    \item Can not exceed a three-month duration even if consideration is given.
\end{enumerate}
\item A criminal actor committed the crime of solicitation. Which of the following does not apply to solicitation?
\begin{enumerate}[label=(\alph*)]
    \item It is an agreement between two or more to commit a crime.
    \item It is solicitation of another to commit a felony.
    \item No act is needed other than the solicitation.
    \item Withdrawal is generally not a defense.
\end{enumerate}
\item In which instance would a state, under the enabling clause of the Fourteenth Amendment, be most able to regulate?
\begin{enumerate}[label=(\alph*)]
    \item A private individual from discriminating against a person based on race.
    \item A private individual from discriminating against a person based on nationality.
    \item A state official from discriminating against a person based on race.
    \item A federal official from discriminating against a person based on nationality.
\end{enumerate}
\end{enumerate}

\section*{True / False (10 marks)}
\textit{Answer True or False}
\begin{enumerate}[label=\arabic*.]
\setcounter{enumi}{5}
\item True or False: In most states, homicide degrees are distinguished by the defendant's state of mind at the time of the killing.
\item True or False: The purpose of the contract must be clearly stated for frustration of purpose to apply.
\item True or False: A warranty of fitness for a particular purpose arises when a seller knows the buyer's specific needs and recommends a product accordingly.
\item True or False: Executory interests provide an immediate right to possession upon the occurrence of a specified event, similar to remainders.
\item True or False: An attorney-client relationship can create an implied-in-fact contract regarding fees when services are accepted without a formal agreement.
\end{enumerate}

\section*{Long Form Response (20 marks)}
\textit{Answer each question with a clear and complete explanation.}
\begin{enumerate}[label=\arabic*.]
\setcounter{enumi}{10}
\item Discuss the implications of the President's power to appoint officials with the Senate's advice and consent, analyzing its impact on the balance of power within the federal government.
\item Discuss the legal implications for a buyer who encounters an unfit, defective, and unsafe product from a merchant, specifically focusing on the warranty of merchantability and its application.
\end{enumerate}

\vfill
\noindent\textit{End of Paper}
\end{document}